\documentclass[11pt, aspectratio=169]{beamer}
\usepackage[style=authoryear-icomp, backend=biber]{biblatex}
\usetheme{jcar}

%\includeonlyframes{a} % Para compilar solo en la que se está trabajando

%\addbibresource{refs.bib}

\begin{document}
\author[Juan C. Arboleda Rivera]{Juan Camilo Arboleda Rivera}
\title{Bootstrap}
%\subtitle{}
%\titlegraphic{\includegraphics[width=2cm]{udea}}
\institute[U. de A.]{Instituto de Biología\\
  Facultad de Ciencias Exactas y Naturales\\
  Universidad de Antioquia\\
  juan.arboleda2@udea.edu.co}
\date{v: \today}
%\subject{}

\begin{frame}[plain]
  \maketitle
\end{frame}

\begin{frame}
 \frametitle{Contenidos}
 \tableofcontents
\end{frame}

\begin{frame}
%[label=a]
{Antes de comenzar...}{}
 \centering
 El bootstrap es una técnica muy usada en Biología\\
 \includegraphics[height=.65\textheight]{Coulson_2006_gadine_phylogeny.jpg}\\
 Por ejemplo, los números en esta filogenia son ``valores de bootstrap'',
 luego explicaremos en qué consisten
\end{frame}

\section{Error estándar e intervalos de confianza}

\begin{frame}
%[label=a]
{}{}
 \centering \Large

 \textbf{Error estándar}\\~\\

 El \textbf{error estándar} (SE) es la desviación estándar de la
 distribución muestral de un estadístico muestral, \textbf{por ejemplo} la
 media\\~\\

 Podemos hablar de error estándar de la desviación estándar, o el error
 estándar de una proporción
\end{frame}

\begin{frame}
%[label=a]
{}{}
 \centering
 \includegraphics[height=.85\textheight]{población-muestras-tortugas.png}\\
 A partir de cada muestra calculamos un estadístico, por ejemplo, la
 proporción de hembras ($ P_h $).
\end{frame}

\begin{frame}
%[label=a]
{}{}
 \centering
 \includegraphics[width=\textwidth]{dist-muestral-proporción.jpg}\\
 Calculamos el error estándar de estos valores, y con este error estándar
 podemos construir un intervalo de confianza del 95\%:\\

  $ IC_{95\%}= P_h \pm \approx 2 \times \text{Error estándar} $
\end{frame}

\begin{frame}
% [label=a]
 {}{}
 \centering \large
 Cuando calculábamos intervalos de confianza \textbf{para la media},
 podíamos hallar el error estándar de la media gracias a que conocíamos la
 ecuación

 \[ SE_{\bar{x}} = \dfrac{s}{\sqrt{n}} \]

 para un estadístico desconocido tendríamos que calcular la desviación
 estándar de muchas muestras, pero...
\end{frame}

\begin{frame}
%[label=a]
{}{}
 \centering \Large
 En la vida real...\\~\\
 ...normalmente solo tenemos \textbf{una} muestra!\\~\\

 \pause

 ¿Cómo sabemos qué tanto varía un estadístico (por ejemplo la proporción de
 hembras), si solo tenemos una muestra?\\~\\

 \pause

 {\LARGE ¡Bootstrap!}
\end{frame}

\section{Técnica de Bootstrap}

\begin{frame}
%[label=a]
{}{}
 \centering \Large
 \includegraphics[height=.85\textheight]{muestra-tortugas-por-sexo.png}\\
 \vspace{-.5em}
 $ \hat{P_h} = 0.7 $\\
 ¿Cuál será el porcentaje ``real''?
\end{frame}

\begin{frame}
%[label=a]
{Creando una ``Población'' ficticia}{}
 \centering \Large
 Imaginemos que la ``Población'' está compuesta de muchas, muchas copias de
 la muestra original...
\end{frame}

\begin{frame}
%[label=a]
{Muestreo con reemplazo}{}
 \centering \Large
 \begin{itemize}[<+->] \itemsep 1em
 \item Para simular una distribución muestral, podemos simplemente tomar
 repetidamente muestras aleatorias de esta ``Población'' hecha de muchas
 copias de nuestra muestra

 \item En la práctica no se pueden tener infinitas copias de la muestra...

 \item ...pero podemos simular esto realizando un \emph{muestreo con
 reemplazo} a partir de la muestra que tenemos (cada individuo puede ser
 seleccionado más de una vez)
 \end{itemize}
\end{frame}

\begin{frame}
%[label=a]
{Conceptos importantes}{}
 \centering \large

 \begin{block}{Muestra de bootstrap}
  Muestra con reemplazo de la muestra original, usando el mismo tamaño
  muestral
 \end{block}

 \includegraphics<1>[width=.9\linewidth]{bootstrap-sample.png}

 \begin{block}<2->{Estadístico de bootstrap}
  Es un estadístico calculado a partir de una muestra de bootstrap
 \end{block}

 \begin{block}<3>{Distribución de bootstrap}
  Es la distribución de muchos estadísticos de bootstrap
 \end{block}
\end{frame}

\begin{frame}
%[label=a]
{}{}
 \centering
 \includegraphics[width=\linewidth]{bootstrap-dist.png}
\end{frame}

\begin{frame}
%[label=a]
{}{}
 \centering
 \includegraphics[height=\textheight]{intro-bootstrap.jpeg}
\end{frame}

\begin{frame}
%[label=a]
{Muestra de bootstrap}{}
% \centering
 \Large
 Supongamos que la muestra original tiene los siguientes datos:
 \[ 7, 8, 9, 42, 42 \]

 ¿Puede la siguiente ser una muestra de bootstrap?
 \only<1,2>{ \color{azul} \textbf{
  \[ 9, 42, 8, 7, 11  \]
  \pause[2]
  No, porque $ 11 $ no es un valor presente en la muestra original
 }}
 \only<3,4>{ \color{azul} \textbf{
  \[ 7, 8, 9, 42 \]
  \pause[4]
  No, porque una muestra de bootstrap debe tener el mismo tamaño que la
  muestra original
 }}
 \only<5,6>{ \color{azul} \textbf{
  \[ 9, 42, 8, 7, 9  \]
  \pause[6]
  Sí, porque pudo haber sido obtenida mediante muestreo con reemplazo de la
  muestra original y es del mismo tamaño
  }}
\end{frame}

\begin{frame}
%[label=a]
{Muestra de bootstrap}{}
 \Large
 Tienes una muestra de tamaño $ n = 50 $. Haces un muestreo con reemplazo
 1000 veces para obtener 1000 muestras de bootstrap\\~\\

 \only<1>{
  ¿Cuál es el tamaño de muestra de cada muestra de bootstrap?\\~\\

  A) 50\\
  B) 1000
 }
 \only<2>{
   ¿Cuántos estadísticos de bootstrap obtendrás?\\~\\

   A) 50\\
   B) 1000
  }
\end{frame}

\begin{frame}
%[label=a]
{¿De dónde viene el nombre Bootstrap?}{}
 \centering
 \begin{columns}
  \column{.25\linewidth}
   \includegraphics[width=\linewidth]{bootstraps1.jpg}
  \column{.6\linewidth}
   \includegraphics[width=\linewidth]{bootstraps2.jpg}
 \end{columns}
 \vspace{1em}
 «Pull oneself up by one's own bootstraps»:\\Levantarse a sí mismo con las
 orejas de las botas\\
 Metáfora para llevar a cabo una tarea ``imposible'' sin ayuda externa
\end{frame}

\begin{frame}
%[label=a]
{}{}
 \footnotesize

 {\centering
 \includegraphics[height=.6\textheight]{bootstrap-phylogenetic-trees.png}\\}
 \begin{enumerate}
 \item Se alinean las secuencias homólogas de ADN de las especies estudiadas.

 \item Se construyen conjuntos de datos Bootstrap mediante el muestreo
 aleatorio con reemplazo de las columnas del alineamiento. Cada conjunto
 conserva el mismo número de posiciones que el alineamiento original.

 \item Se construye un árbol filogenético para cada conjunto.

 \item Se comparan los árboles obtenidos y se calcula la frecuencia con la
 que aparece cada agrupación de especies.
 \end{enumerate}
\end{frame}

\begin{frame}
% [label=a]
 {}{}
 \centering
 \includegraphics[height=.9\textheight]{Coulson_2006_gadine_phylogeny.jpg}
\end{frame}

\section{Intervalos de confianza por bootstrap}

\begin{frame}
%[label=a]
{Error estándar y distribución de bootstrap}{}
 \centering \Large

 \begin{itemize}[<+->] \itemsep 1em
 \item La variabilidad de los estadísticos de bootstrap es similar a la
 variabilidad de los estadísticos muestrales

 \item ¡El \textbf{error estándar} de un estadístico puede ser estimado
 usando la desviación estándar de la distribución de bootstrap!
 \end{itemize}
\end{frame}

\begin{frame}
%[label=a]
{La magia del bootstrapping}{}
 \centering \Large

 \begin{itemize} \itemsep 1em
 \item Podemos usar bootstrapping para estimar la incertidumbre o error
 estándar de muchos estadísticos muestrales

 \item Si tenemos datos de una muestra, podemos usar bootstrapping para
 crear un intervalo de confianza del 95\% para casi cualquier parámetro
 \end{itemize}
\end{frame}

\begin{frame}
%[label=a]
{Intervalos de confianza mediante Bootstrap}{}
 \centering \Large

 \begin{description}
 \item[Opción 1:] Estimar el error estándar del estadístico calculando la
 desviación estándar de la distribución de bootstrap, y luego generar un
 intervalo de confianza del 95\% calculando:

 \[ IC_{95\%}= \text{estadístico} \pm 2 \cdot SE \]

 \item[Opción 2:] Generar un intervalo de confianza del X\% como el rango
 para los X\% valores centrales de los estadísticos de bootstrap
 \end{description}
\end{frame}

\begin{frame}
%[label=a]
{Método de los percentiles para intervalos de confianza}{}
 \centering \large

 \begin{itemize} \itemsep 1em
 \item Para un intervalo de confianza del X\%, mantener los X\% valores
 centrales de los estadísticos de bootstrap

 \item Por ejemplo, para un intervalo de confianza del 99\%, mantener el
 99\% de los valores centrales, dejando 0.05\% en cada cola

 \item El intervalo de confianza sería:\\
 \[ IC_{99\%} = [\text{Percentil 0.05}, \text{Percentil 99.5}] \]\\
 donde estos percentiles se refieren a la distribución de bootstrap
 \end{itemize}
\end{frame}

\begin{frame}
%[label=a]
{Dos métodos}{}
 \centering \Large

 \begin{itemize} \itemsep 1em
 \item Para una distribución de bootstrap simétrica, con forma de campana,
 usar el método del error estándar o el método de los percentiles dará
 intervalos de confianza del 95\% muy similares

 \item Si la distribución de bootstrap no tiene forma de campana o si se
 quiere un nivel de confianza diferente del 95\%, use el método de los
 percentiles
 \end{itemize}
\end{frame}

\begin{frame}
%[label=a]
{Criterios para construcción de Intervalos de Confianza mediante Bootstrap}{}
 \Large

 Usar el método del percentil para un intervalo de confianza mediante
 bootstrap funciona para cualquier estadístico, siempre y cuando la
 distribución de bootstrap sea:

 \begin{itemize}
  \item Aproximadamente continua
 \end{itemize}

 \vspace{1em}
 Usar el método del error estándar requiere también que la distribución sea:

 \begin{itemize}
  \item Aproximadamente simétrica
  \item Con forma de campana
 \end{itemize}
\end{frame}

\begin{frame}
%[label=a]
{Número de muestras de bootstrap}{}
 \centering \large

 \begin{itemize} \itemsep 1em
 \item Cuando se usa bootstrap, se puede obtener un intervalo de confianza
 ligeramente diferente cada vez. Esto es normal.

 \item Entre mayor cantidad de muestras de bootstrap se usen, más precisa
 será la respuesta

 \item En muchos casos de la vida real se usan 10.000 o hasta 100.000
 muestras de bootstrap
 \end{itemize}
\end{frame}

\begin{frame}
%[label=a]
{Resumen}{}
 \centering \large

 \begin{itemize} \itemsep 1em
 \item El error estándar de un estadístico es la desviación estándar del
 estadístico entre muchas muestras, el cual puede ser estimado a partir de
 una distribución de bootstrap

 \item Se pueden crear intervalos de confianza usando el error estándar o
 los percentiles de una distribución de bootstrap
 \end{itemize}
\end{frame}

\end{document}
